% !TeX program = xelatex
% ============================================================
% vkr/vkr.tex — Выпускная квалификационная работа
% ОП «Программная инженерия» ФКН НИУ ВШЭ.
% Оформление: ГОСТ 7.32-2017 «Отчёт о научно-исследовательской
% работе. Структура и правила оформления». Класс extreport — главы.
%
% Структура (методичка ОП ПИ + ГОСТ 7.32):
%   Титульный лист → Реферат → Abstract → Содержание →
%   Термины и определения → Список сокращений → Введение →
%   3 главы → Заключение → Список источников → Приложения.
%
% Как пользоваться: пишите свой текст вместо пояснений (строки,
% начинающиеся с «%», в PDF не попадают). Данные о себе, научном
% руководителе, группе и годе задайте в настройках проекта —
% они подставятся в титул и реферат автоматически (макросы \hse…).
% ============================================================
\documentclass[12pt,a4paper]{extreport}
\input{preamble}
\begin{document}
\input{title}

% ============================================================
\chapter*{Реферат}
\addcontentsline{toc}{chapter}{Реферат}
% ============================================================
% Реферат (ГОСТ 7.32-2017 §5.3) — краткое и точное изложение работы.
% Структура одного абзаца (150–250 слов): проблема и актуальность →
% цель → что сделано (требования, проектирование, реализация) →
% главный результат и область его применения.

Здесь приводится краткое изложение работы: проблема и её
актуальность, цель работы, выполненные задачи
((* if project.kind == 'research' *))(аналитический обзор состояния проблемы, разработка
модели и метода решения, экспериментальная проверка результатов)((* else *))(анализ предметной
области, формирование требований, проектирование и реализация)((* endif *)),
а также главный полученный результат и область его применения.

\medskip
\textbf{Работа содержит} \pageref{LastPage}~страниц, \ldots~глав,
\ldots~рисунков, \ldots~таблиц, \ldots~источников и \ldots~приложения.

\medskip
% Ключевые слова (ГОСТ 7.32-2017 §5.3.2.1, §6.12.2): от 5 до 15 слов или
% словосочетаний, ПРОПИСНЫМИ буквами, в именительном падеже, в строку
% через запятую, без точки в конце.
\textbf{Ключевые слова:} ПРОГРАММНАЯ ИНЖЕНЕРИЯ, \ldots

\clearpage

% ============================================================
\chapter*{Abstract}
\addcontentsline{toc}{chapter}{Abstract}
% ============================================================
% English abstract (150–250 words) — обязателен для ВКР ФКН НИУ ВШЭ
% даже если работа написана на русском. Это перевод реферата.

This section contains a concise English summary of the thesis: the
problem and its relevance, the objective, the work carried out
((* if project.kind == 'research' *))(a systematic review of the state of the art, the proposed
model and method, and experimental evaluation)((* else *))(domain analysis, requirements,
design and implementation)((* endif *)), and the
main result with its area of application.

\medskip
\textbf{The work contains} \pageref{LastPage}~pages, \ldots~chapters,
\ldots~figures, \ldots~tables, \ldots~sources and \ldots~appendices.

\medskip
\textbf{Keywords:} SOFTWARE ENGINEERING, \ldots

\clearpage

% ── Содержание / Contents ───────────────────────────────────
\tableofcontents
\clearpage


% ============================================================
\chapter*{Термины и определения}
\addcontentsline{toc}{chapter}{Термины и определения}
% ============================================================
% Глоссарий ключевых терминов работы (description-список). Удалите
% раздел, если специальная терминология не используется.

В настоящей работе применяются следующие термины с соответствующими
определениями.

\begin{description}[leftmargin=0pt, labelindent=0pt, font=\bfseries, style=nextline]
  \item[API] прикладной программный интерфейс — набор соглашений
    для взаимодействия программных компонентов.
  \item[Прототип] частичная реализация системы, создаваемая для
    проверки ключевых архитектурных и функциональных решений.
  \item[\ldots] \ldots\ % добавьте остальные термины
\end{description}

\clearpage

% ============================================================
\chapter*{Список сокращений и условных обозначений}
\addcontentsline{toc}{chapter}{Список сокращений и условных обозначений}
% ============================================================
% Перечень аббревиатур и условных обозначений (ГОСТ 7.32-2017 §5.3).
% Можно оформить description-списком (как ниже) либо таблицей.

В настоящей работе используются следующие сокращения и условные
обозначения.

\begin{description}[leftmargin=0pt, labelindent=0pt, font=\bfseries, style=nextline]
  \item[API] Application Programming Interface (программный интерфейс приложения).
  \item[ГОСТ] государственный стандарт.
  \item[СУБД] система управления базами данных.
  \item[\ldots] \ldots\ % добавьте остальные сокращения
\end{description}

\clearpage

% ============================================================
\chapter*{Введение}
\addcontentsline{toc}{chapter}{Введение}
% ============================================================
% Введение (ГОСТ 7.32-2017 §5.4, методичка ОП ПИ) обязательно
% раскрывает: актуальность, цель, задачи, объект, предмет, методы,
% новизну и практическую значимость, структуру работы.
% Рекомендуемый объём — не менее 500 слов (2–4 страницы).

\section*{Актуальность темы}
Обоснуйте, почему задача важна сейчас: какая проблема решается, кого
она затрагивает и почему существующие подходы недостаточны. Ссылку на
источник оформляйте командой \verb|\cite|, например~\cite{article}.

\section*{Цель и задачи работы}
\textbf{Цель работы} — \ldots\ % одно предложение: что именно создаётся

\textbf{Задачи:}
\begin{enumerate}[label=\arabic*), leftmargin=1.25cm]
((* if project.kind == 'research' *))
  \item провести аналитический обзор литературы и существующих подходов;
  \item выявить нерешённые проблемы и сформулировать задачу исследования;
  \item разработать модель (метод, алгоритм) решения поставленной задачи;
  \item провести вычислительный эксперимент и проанализировать результаты.
((* else *))
  \item проанализировать предметную область и существующие решения;
  \item сформировать функциональные и нефункциональные требования;
  \item спроектировать архитектуру и модель данных;
  \item реализовать прототип и провести его тестирование.
((* endif *))
\end{enumerate}

\section*{Объект и предмет}
\textbf{Объект исследования} — \ldots

\textbf{Предмет исследования} — \ldots

((* if project.kind == 'research' *))
\section*{Методы исследования}
Перечислите методы исследования (аналитический обзор, формальное
моделирование, статистический анализ, вычислительный эксперимент) и
инструментальные средства, использованные в работе.
((* else *))
\section*{Методы и инструменты}
Перечислите методы (анализ литературы, сравнение аналогов,
моделирование, эксперимент) и инструменты, использованные в работе.
((* endif *))

\section*{Научная новизна и практическая значимость}
Опишите, что нового предложено в работе и где результат может быть
применён на практике.

\section*{Структура работы}
Работа состоит из введения, основной части, заключения, списка
использованных источников и приложений. В первой главе \ldots; во
второй главе \ldots; в третьей главе \ldots\ Оформление работы
выполнено по ГОСТ~7.32-2017~\cite{gost732}.

\clearpage

((* if project.kind == 'research' *))
% ============================================================
\chapter{ОБЗОР ЛИТЕРАТУРЫ И ПОСТАНОВКА ЗАДАЧИ}
% ============================================================

\section{Современное состояние проблемы}
Дайте аналитический обзор предметной области: ключевые понятия,
существующие теоретические результаты и подходы к решению
рассматриваемой проблемы. Ссылайтесь на источники командой
\verb|\cite|, например~\cite{article}.

\section{Сравнительный анализ существующих подходов}
Сопоставьте известные методы (подходы, модели) по значимым для
исследования критериям; результат сравнения приведён в
таблице~\ref{tab:approaches}. В ячейках используйте \verb|\cmpplus| (свойство
присутствует), \verb|\cmpminus| (отсутствует), \verb|\cmpplanned| (частично).

\begin{table}[H]
\caption{Сравнение существующих подходов}
\label{tab:approaches}
\centering
\begin{tabular}{|>{\raggedright\arraybackslash}p{5cm}|c|c|c|}
\hline
\textbf{Критерий} & \textbf{Подход~А} & \textbf{Подход~Б} & \textbf{Предлагаемый} \\
\hline
Критерий~1 & \cmpplus  & \cmpminus  & \cmpplus \\
\hline
Критерий~2 & \cmpminus & \cmpplanned & \cmpplus \\
\hline
\end{tabular}
\end{table}

\section{Нерешённые проблемы и постановка задачи исследования}
На основе обзора выделите пробелы в существующих подходах,
сформулируйте гипотезу исследования и дайте формальную постановку
решаемой задачи (входные данные, искомый результат, ограничения,
критерии качества).

\section{Выводы по главе}
Кратко сформулируйте итоги обзора и обоснуйте выбранное направление
исследования.

\clearpage

% ============================================================
\chapter{МОДЕЛЬ И МЕТОД РЕШЕНИЯ}
% ============================================================

\section{Формальная постановка}
Введите обозначения и строго сформулируйте исследуемую задачу: входные
данные, искомый результат, ограничения и критерий качества (например,
целевую функцию $f(x)\to\min$ или проверяемую гипотезу).

\section{Предлагаемая модель (метод, алгоритм)}
Опишите предлагаемый метод (модель, алгоритм) решения задачи: основные
шаги, используемые допущения и параметры. Формальные преобразования
поясняйте по ходу изложения.

\section{Обоснование корректности подхода}
Приведите теоретическое обоснование предлагаемого метода: доказательство
свойств, оценки вычислительной сложности или точности, либо
содержательную аргументацию.

\section{Выводы по главе}
Сформулируйте ключевые теоретические результаты главы.

\clearpage

% ============================================================
% Третья глава исследовательской ВКР зависит от характера работы
% (арка академической ВКР ОП ПИ НЕ фиксирована):
%   • экспериментальное исследование — оставьте «ЭКСПЕРИМЕНТАЛЬНОЕ
%     ИССЛЕДОВАНИЕ» (методика, данные, метрики, результаты — как ниже);
%   • инженерная/теоретическая работа без вычислительного эксперимента —
%     замените главу на «РЕАЛИЗАЦИЯ И АПРОБАЦИЯ» (средства реализации,
%     программная реализация, тестирование, апробация). Главы-эксперимента
%     в таком случае может не быть.
% ============================================================
\chapter{ЭКСПЕРИМЕНТАЛЬНОЕ ИССЛЕДОВАНИЕ}
% ============================================================

\section{Методика эксперимента}
Опишите план эксперимента: исследовательские вопросы, условия и
процедуру, в которых проверяется гипотеза.

\section{Данные и метрики}
Опишите используемые наборы данных (объём, источник, подготовка) и
метрики оценки качества результатов.

\section{Результаты и их анализ}
Приведите полученные результаты (таблица~\ref{tab:results}), сравните их
с базовыми методами и оцените значимость различий.

\begin{table}[H]
\caption{Результаты эксперимента}
\label{tab:results}
\centering
\begin{tabular}{|>{\raggedright\arraybackslash}p{5cm}|c|c|}
\hline
\textbf{Метод} & \textbf{Метрика~1} & \textbf{Метрика~2} \\
\hline
Базовый метод      & \ldots & \ldots \\
\hline
Предлагаемый метод & \ldots & \ldots \\
\hline
\end{tabular}
\end{table}

\section{Обсуждение}
Проинтерпретируйте результаты, укажите ограничения исследования и
условия применимости полученных выводов.

\section{Выводы по главе}
Сформулируйте итоги экспериментального исследования.

\clearpage

((* else *))
% ============================================================
\chapter{АНАЛИЗ ПРЕДМЕТНОЙ ОБЛАСТИ}
% ============================================================

\section{Анализ предметной области}
Опишите предметную область, ключевые понятия и проблему, которую
решает работа.

\section{Обзор существующих решений}
Сравните существующие аналоги по значимым критериям. На каждую
таблицу должна быть ссылка в тексте (например: сравнение приведено в
таблице~\ref{tab:analogues}). В ячейках используйте \verb|\cmpplus| (есть),
\verb|\cmpminus| (нет), \verb|\cmpplanned| (частично/в планах).

\begin{table}[H]
\caption{Сравнение существующих решений}
\label{tab:analogues}
\centering
\begin{tabular}{|>{\raggedright\arraybackslash}p{5cm}|c|c|c|}
\hline
\textbf{Критерий} & \textbf{Аналог~А} & \textbf{Аналог~Б} & \textbf{Предлагаемое} \\
\hline
Критерий~1 & \cmpplus  & \cmpminus  & \cmpplus \\
\hline
Критерий~2 & \cmpminus & \cmpplanned & \cmpplus \\
\hline
\end{tabular}
\end{table}

\section{Требования к системе}
Сформулируйте функциональные и нефункциональные требования. Если
требования определены в Техническом задании макросом
\verb|\req{ТЗ-Ф-01}{...}|, ссылайтесь на них в тексте через
\verb|\reqref{ТЗ-Ф-01}| — панель «Трассировка требований» свяжет их
автоматически. Перечень функциональных требований приведён в
таблице~\ref{tab:func-req}.

\begin{table}[H]
\caption{Функциональные требования}
\label{tab:func-req}
\centering
\begin{tabular}{|c|>{\raggedright\arraybackslash}p{11cm}|}
\hline
\textbf{Идентификатор} & \textbf{Описание} \\
\hline
ТЗ-Ф-01 & \ldots \\
\hline
\end{tabular}
\end{table}

\section{Выводы по главе}
Кратко сформулируйте итоги анализа и обоснуйте постановку задачи.

\clearpage

% ============================================================
\chapter{ПРОЕКТИРОВАНИЕ}
% ============================================================

\section{Архитектура решения}
Опишите общую архитектуру. Каждый рисунок должен быть упомянут в
тексте (например: <<общая схема приведена на рисунке~\verb|\ref{fig:arch}|>>).
Иллюстрацию можно вставить как \verb|\includegraphics| (отсутствующий
файл заменяется заглушкой) либо нарисовать средствами \texttt{tikz}.

\section{Модель данных}
Опишите основные сущности и связи (ER-диаграмма или текстовое
описание).

\section{Выбор технологического стека}
Перечислите языки, фреймворки и библиотеки с обоснованием выбора.

\section{Выводы по главе}
Сформулируйте принятые проектные решения.

\clearpage

% ============================================================
\chapter{РЕАЛИЗАЦИЯ И ТЕСТИРОВАНИЕ}
% ============================================================

\section{Детали реализации}
Опишите ключевые модули и нетривиальные решения. Фрагменты кода
оформляйте окружением \texttt{lstlisting} (см. листинг~\ref{lst:example}).

\begin{lstlisting}[language=Python, caption={Пример функции}, label={lst:example}]
def greet(name: str) -> str:
    return f"Hello, {name}!"
\end{lstlisting}

\section{Тестирование}
Опишите стратегию тестирования (модульные, интеграционные, e2e) и
приведите результаты в таблице~\ref{tab:test-results}.

\begin{table}[H]
\caption{Результаты тестирования}
\label{tab:test-results}
\centering
\begin{tabular}{|>{\raggedright\arraybackslash}p{4cm}|>{\raggedright\arraybackslash}p{8cm}|c|}
\hline
\textbf{Вид тестирования} & \textbf{Что проверяется} & \textbf{Результат} \\
\hline
Модульное & \ldots & \ldots \\
\hline
Интеграционное & \ldots & \ldots \\
\hline
\end{tabular}
\end{table}

\section{Апробация и результаты}
Опишите, где и как решение проверено, и какие результаты получены.

\section{Выводы по главе}
Сформулируйте итоги реализации и тестирования.

\clearpage
((* endif *))

% ============================================================
\chapter*{Заключение}
\addcontentsline{toc}{chapter}{Заключение}
% ============================================================
% Заключение (ГОСТ 7.32-2017 §5.5): краткие итоги по каждой задаче
% из введения, главный результат, соответствие цели и перспективы
% дальнейшего развития.

По результатам выполнения работы:
\begin{enumerate}[label=\arabic*), leftmargin=1.25cm]
  \item \ldots\ % итог по задаче 1
  \item \ldots\ % итог по задаче 2
\end{enumerate}

Цель работы достигнута. Перспективы дальнейшего развития: \ldots

\clearpage

% ============================================================
% Список использованных источников
% ============================================================
% Список источников по ГОСТ Р 7.0.5-2008 (ссылки) + ГОСТ 7.1-2003
% (описание) + ГОСТ 7.82-2001 (электронные ресурсы) — методические
% указания ФКН / ОП ПИ (НЕ ГОСТ Р 7.0.100-2018). Единый алфавитный
% нумерованный список (сначала латиница, затем кириллица), разделитель
% областей «. —». На каждый источник должна быть хотя бы одна ссылка
% \cite{} в тексте. По методичке ОП ПИ список ВКР содержит не менее
% 20 источников (для исследовательской работы — преимущественно научные).
\phantomsection
\addcontentsline{toc}{chapter}{Список использованных источников}
\begin{thebibliography}{99}

  % Книга (ГОСТ 7.1-2003)
  \bibitem{book}
    Фамилия, И.~О. Заглавие книги : сведения, относящиеся к заглавию /
    И.~О.~Фамилия. — М.: Издательство, 2023. — 320~с.

  % Статья из журнала (ГОСТ Р 7.0.5-2008 / 7.1-2003)
  \bibitem{article}
    Фамилия, И.~О. Название статьи / И.~О.~Фамилия, А.~Б.~Соавтор //
    Название журнала. — 2024. — № 1. — С.~10--20.

  % Электронный ресурс (ГОСТ 7.82-2001)
  \bibitem{webresource}
    Название ресурса [Электронный ресурс]. — Режим доступа:
    \url{https://example.org}, свободный (дата обращения: 01.01.2026).

  % Нормативный документ
  \bibitem{gost732}
    ГОСТ 7.32--2017. Система стандартов по информации, библиотечному
    и издательскому делу. Отчёт о научно-исследовательской работе.
    Структура и правила оформления. — М.: Стандартинформ, 2017. — 27~с.

\end{thebibliography}

\clearpage

% ============================================================
% Приложения (ГОСТ 7.32-2017 §6.14)
% ============================================================
% Приложения обозначаются русскими буквами А, Б, В… На каждое
% приложение должна быть ссылка в тексте (например: см.
% приложение~\ref{app:listing}).
\appendix
\renewcommand{\thechapter}{\Asbuk{chapter}}
\titleformat{\chapter}[hang]
  {\normalfont\large\bfseries\centering}{Приложение~\thechapter}{1ex}{}

\chapter{Листинги исходного кода}
\label{app:listing}
Разместите здесь объёмные листинги, схемы, таблицы данных и другие
вспомогательные материалы.


\end{document}
